\section{Conclusions}
\label{sec:conclusion}

% PAGE TARGET: ~2 pages.
% Day 6 (5/30).
% CUED wants: main findings + future work + REFLECTION on what worked,
% what did not, how I could have improved.

\subsection{Summary of Findings}
% Restate the three contributions from Section 1.4, now backed by
% concrete numbers from Section 4:
%   - Pipeline reaches sub-1\% density error on calibrated folios.
%   - Synthetic phantom characterisation confirms SNR-invariance and
%     period accuracy in 10--50 px range.
%   - Failure modes localised (octave aliasing, polarity).

\subsection{Reflection on the Project}
% CUED-required reflection paragraph (a few hundred words):
% - What worked well: staged modular pipeline made adding new detectors
%   cheap; phantom framework caught the projection-autocorrelation bug
%   early.
% - What did not: spreadsheet GT mismatch was discovered late -- earlier
%   manual GT collection would have shifted Results emphasis.
% - With hindsight: I would have built the phantom sweep before the real
%   detector, not after; would have set up manual GT in Lent term.

\subsection{Future Work}
% - Automatic polarity detection (drop wire\_is\_darker).
% - Joint chain-line detection: currently only laid lines.
% - Mould-mate matching across folios using per-folio fingerprint
%   (period + FWHM distribution + chain-line spacing).
% - Confidence-aware aggregation across multi-page codices.
